% ══════════════════════════════════════════════════════════════════════════════
\section{Literature Review}
\label{sec:litreview}
% ══════════════════════════════════════════════════════════════════════════════

\subsection{Climate, Income Shocks, and Conflict}

A large and growing empirical literature examines the relationship between environmental shocks, economic conditions, and violent conflict, with Sub-Saharan Africa as the primary focus. The foundational contribution of \citet{MiguelSatyanathSergenti2004} establishes that negative rainfall shocks significantly increase the probability of civil conflict onset in Sub-Saharan African countries, a result widely interpreted through the opportunity-cost channel in largely agrarian economies where rainfall is the primary determinant of agricultural income. This study became the methodological template for the field: using rainfall variation as a source of plausibly exogenous income shocks, it demonstrates that economic deterioration increases the risk of large-scale political violence.

Subsequent work has extended this approach across broader samples, alternative climatic dimensions, and different conflict definitions. \citet{DellJonesOlken2014} document that higher temperatures reduce economic output and increase political instability, with effects concentrated in low-income countries. \citet{BurkeHsiangMiguel2015} conduct a meta-analysis of 55 studies and find a robust statistical association between climatic anomalies---particularly temperature extremes and rainfall deficits---and a wide range of conflict outcomes including interpersonal violence, political conflict, and institutional breakdown. A recent comprehensive synthesis by \citet{BurkeFergusonHsiangMiguel2024} reviews the expanded literature and concludes that the empirical association between climate and conflict is robust, but that effect magnitudes vary considerably across contexts, and that understanding the economic transmission mechanisms---especially agriculture---is essential for interpreting reduced-form estimates.

Despite methodological sophistication, a key limitation of this literature is the treatment of climatic variables as proxies for economic welfare without modelling the agricultural transmission mechanism explicitly. A given rainfall deviation translates into very different yield outcomes depending on crop type, soil conditions, irrigation infrastructure, and local adaptation practices. Studies that rely on rainfall as an instrument for agricultural income effectively estimate local average treatment effects that may conflate the direct physiological and psychological effects of climate on behaviour with the indirect effects operating through production and income.

\subsection{Agriculture as the Transmission Mechanism}

A more targeted strand of the literature focuses explicitly on agricultural productivity as the economic channel through which climate affects conflict. \citet{CoutteynierSoubeyran2014} show, using a panel of African countries, that drought-induced agricultural losses---proxied by the Palmer Drought Severity Index---significantly increase civil conflict incidence. \citet{HarariLaFerrara2018} extend this approach to a spatially disaggregated analysis, demonstrating that within-country conflict incidence is elevated in growing-season areas affected by negative vegetation anomalies, and that this relationship holds at fine spatial scales consistent with localised agricultural stress rather than country-level income effects.

Critically, the sign of the relationship between agricultural conditions and conflict is not uniformly negative across this literature, reflecting the competing mechanisms discussed in Section~\ref{sec:theory}. \citet{DubeVargas2013} provide particularly clean evidence on this point using Colombian commodity price shocks. They show that an increase in the international price of oil---a capital-intensive commodity that is easily appropriable by armed groups---significantly increases conflict in oil-producing regions, while an increase in the price of coffee---a labour-intensive crop that raises rural wages and hence opportunity costs---reduces conflict. This asymmetry demonstrates that the direction of the agricultural income--conflict relationship depends on the appropriability of the commodity and the relative magnitudes of the opportunity-cost and predation channels.

Evidence from Sub-Saharan Africa is similarly mixed. \citet{Berman_etal2017} examine the relationship between commodity price movements and conflict across a panel of African countries, finding that positive price shocks to resources controlled by armed groups increase conflict while shocks to broadly distributed agricultural commodities do not. \citet{AldrichSarr2020} find that food price increases in urban areas generate civil unrest, while rural agricultural price increases reduce it---consistent with an opportunity-cost mechanism for producers and a cost-of-living channel for non-producers. This heterogeneity underscores that the relationship between agricultural conditions and conflict is deeply context-dependent, varying with production structure, market integration, and the distribution of gains from agricultural improvement.

\subsection{Mixed Evidence and Contextual Heterogeneity}

The empirical literature as a whole exhibits substantial heterogeneity in estimated effects, reflecting genuine variation in the mechanisms at work across different institutional and agro-ecological settings. Several factors have been identified as important moderators of the agriculture--conflict relationship.

First, institutional quality---particularly state capacity and the rule of law---shapes whether agricultural shocks generate violence or are absorbed through peaceful adjustment mechanisms. Weak states with limited enforcement capacity are more likely to experience predation in response to positive agricultural shocks, while strong institutions can buffer communities against negative shocks through redistribution and public provision \citep{BesleyPersson2011}. Second, crop appropriability determines the relative strength of the predation and opportunity-cost channels: easily stored and transported crops (grains, oil seeds) are more vulnerable to looting and taxation by armed groups than perishable horticultural produce \citep{RapoportTernstrm2011}. Third, the spatial scale of analysis matters: country-level estimates often mask within-country heterogeneity that is revealed only at subnational resolution \citep{HarariLaFerrara2018}.

A further source of heterogeneity is the time horizon over which conflict outcomes are measured. Yield shocks may generate immediate responses through changes in recruitment and opportunity costs, but their effects on organised conflict could take longer to materialise as groups form, negotiate, and mobilise. Studies that measure conflict contemporaneously with the shock may miss delayed effects captured by longer forward windows, motivating the use of 3-, 6-, and 12-month forward windows in the present analysis.

\subsection{Measurement: From Climate Proxies to Predicted Yields}

A persistent limitation of the empirical literature is its reliance on climatic variables rather than direct measures of agricultural productivity. While rainfall and temperature anomalies provide plausibly exogenous variation, they do not map uniformly into crop yields, and their use as proxies conflates multiple transmission channels. This measurement limitation is increasingly recognised as central to interpreting estimated effects \citep{BurkeFergusonHsiangMiguel2024}.

Direct measures of agricultural productivity offer a more economically precise test of the theoretical mechanisms, but are constrained by data availability. Household and farm surveys provide high-quality yield information but are spatially sparse, temporally irregular, and often subject to measurement error. These constraints restrict the ability of researchers to examine subnational variation in the yield--conflict relationship using survey-based measures alone.

Recent advances in remote sensing and machine learning offer a promising alternative. A substantial literature in agricultural science and development economics demonstrates that satellite-derived vegetation indices, combined with ground-truth survey data and machine-learning methods, can accurately predict crop yields at high spatial resolution in data-scarce environments \citep{BurkeLobell2017, LobellEtal2015, YehEtal2020}. Vegetation indices such as NDVI, EVI, and GCVI capture the greenness and vigour of vegetation during the growing season and correlate strongly with photosynthetic activity and biomass accumulation \citep{BurkeLobell2017}. Temperature-based stress indicators including growing-degree days (GDD) and killing-degree days (KDD) provide additional information on crop development and heat stress \citep{Schlenker_Roberts2009}.

The application of prediction-based measurement strategies to economic research has been formalised by \citet{MullainathanSpiess2017} and \citet{AtheyImbens2019}, who argue that machine-learning predictions can serve as economically meaningful measures of otherwise unobservable variables, provided the prediction model is trained on appropriate ground-truth data and the exclusion restriction linking predictions to outcomes operates only through the predicted quantity. In the present context, yield predictions based solely on satellite and weather inputs satisfy this exclusion restriction as long as conflict does not affect the satellite-measured vegetation signal at the training locations---a reasonable assumption given that conflict intensity is sparse and localised relative to the broad spatial footprint of satellite imagery.

Related work demonstrates that satellite-based measures can successfully recover economic outcomes in developing country contexts where administrative data is limited \citep{JeanEtal2016, BluMenStockEtal2015}. \citet{JeanEtal2016} show that night-time lights predicted from daytime satellite imagery can proxy for consumption and asset wealth at fine spatial scales, providing a model for the satellite-based yield prediction approach used in this paper. \citet{BurkeLobell2017} directly validate satellite-based yield predictions against LSMS survey observations across multiple African countries, finding spatial $R^2$ values in the range of 0.10--0.30 for XGBoost models---consistent with the model performance achieved in this paper.

\subsection{This Paper's Position in the Literature}

This paper contributes to the literature at the intersection of these strands. Relative to the climate--conflict literature, it moves beyond climate proxies to measure the agricultural shock directly, providing a sharper test of the economic transmission mechanism. Relative to the satellite-based measurement literature, it applies yield predictions at scale to a policy-relevant question---the determinants of conflict---rather than treating prediction as an end in itself. Relative to prior subnational studies of agriculture and conflict in Africa, it covers a broader and more heterogeneous set of countries and years, and employs a richer set of conflict outcome measures and estimators.

The approach most closely related to this paper is \citet{HarariLaFerrara2018}, who also use satellite vegetation data to identify agricultural shocks in a subnational African panel. A key difference is that the present paper takes an additional step: rather than using raw vegetation anomalies as the regressor, it trains a supervised model to map vegetation signals to yield outcomes, producing a more economically interpretable measure that directly captures crop productivity rather than general vegetation greenness. This distinction matters because vegetation anomalies reflect a mixture of signals including land-cover change, non-agricultural vegetation, and seasonal variation, while yield predictions condition on the full harvest-cycle lag structure relevant for crop development.
